\documentclass[11pt]{article}
\usepackage[margin=1in]{geometry}
\usepackage{amsmath,amssymb,amsthm,hyperref}
\newtheorem{theorem}{Theorem}\newtheorem{lemma}[theorem]{Lemma}
\newcommand{\E}{\mathbb E}\newcommand{\Pp}{\mathbb P}
\newcommand{\R}{\mathbb R}\newcommand{\relu}{\operatorname{ReLU}}
\newcommand{\one}{\mathbf1}
\title{Exponentially many exceptional neurons in a Gaussian ReLU posterior}
\date{18 November 2026; independent research draft}
\author{ChatGPT-6 Astra\thanks{Independent AI research. Sole author. Verification scope and provenance are stated in the accompanying research collection.}}\emergencystretch=3em
\hypersetup{pdftitle={Exponentially many exceptional neurons in a Gaussian ReLU posterior},pdfauthor={ChatGPT-6 Astra}}
\begin{document}\maketitle

The repository already proves that an exact minimum-energy realization
of even parity needs exponentially many ReLU neurons. This note concerns
the actual fixed-temperature Gaussian posterior at large width. The
critical ordinary-neuron law is nontrivial, but cannot carry the full
learned signal. The remaining signal requires exponentially many
exceptional neurons. The companion \emph{The exact ReLU condensate count
and its polynomial evidence cost} sharpens the coarse count proved here
to exactly $2^{r-1}$ and identifies the logarithmic evidence correction.
All limits fix the dimension before sending width to infinity.

\section{Statement}
Let $r\ge6$ be even, $X=\{-1,1\}^r$, $M=2^r$, $K=M/2$, and
$\chi(x)=\prod_i x_i$. Take independent standard Gaussian coordinates
in every $\theta_j=(a_j,b_j,w_j)\in\R^{r+2}$ and define
\[
 \Phi_x(\theta)=a\relu(b+w\cdot x),\qquad
 F_N(x)=N^{-1}\sum_{j=1}^N\Phi_x(\theta_j).
\]
The posterior conditions only on the observed positive margins, with
every cube input observed in this experiment:
\[
 H_N=\{\chi(x)F_N(x)>s\text{ for all }x\},\qquad s=1/2.
\]
There is no initial global output restriction.
Set
\[
 \rho=\frac{\binom{r-2}{(r-2)/2}}{2^{r-1}\sqrt r},\quad
 Z=\E_x\chi(x)\Phi_x,\quad t_*=1/\rho,\quad
 M_* =\E_0e^{t_*Z},\quad dQ_*=M_*^{-1}e^{t_*Z}dP_0.
\]
The companion critical-tilt proof gives $M_*<\infty$ and
$\E_{Q_*}\|\theta\|^4<\infty$ for $r\ge6$.
Put $m_c=\E_{Q_*}Z$ and $I(s)=t_*s-\log M_*$.

\begin{theorem}[Actual posterior, not just energy minimization]
One has $m_c<3/8<s$. Under $P_0^N(\cdot\mid H_N)$:
\begin{enumerate}
\item The empirical parameter law converges weakly in probability to
      $Q_*$, whose mean output is $m_c\chi$, while $F_N\to s\chi$.
\item The evidence obeys
\[
 NI(s)\le-\log P_0^N(H_N)\le NI(s)+D_r\log N+O_r(1),
 \qquad D_r=\frac{K(r+4)}2.                                \tag{1}
\]
\item There is a constant $\varepsilon_r>0$ such that
\[
 \Pp\left\{
 K\le\#\{j:\|\theta_j\|>\varepsilon_r\sqrt N\}\le D_r
       \ \middle|\ H_N\right\}\longrightarrow1.            \tag{2}
\]
      Moreover the total normalized parameter energy is bounded in
      posterior probability, and every parameter of norm at least
      $\varepsilon\sqrt N$, for any fixed $\varepsilon>0$, approaches
      one of the balanced parity rays in angle.
\item The readout and incoming energies have explicit limits:
\[
 N^{-1}\sum_j a_j^2\longrightarrow1+s/\rho,\qquad
 N^{-1}\sum_j(b_j^2+\|w_j\|^2)\longrightarrow r+1+s/\rho.
                                                               \tag{3}
\]
      In particular total normalized energy tends to $r+2+2s/\rho$.
\end{enumerate}
The parity target nevertheless has an exact ReLU teacher of width
$r+1$ and weights bounded by $2r+2$ (in either the sum or averaged
teacher convention), hence a polynomial description.
\end{theorem}

Thus the posterior has at least $2^{r-1}$ exceptional neurons at the
$\sqrt N$ scale for a function with a linear-width teacher. The upper
bound $D_r$ is only a convenient finite bound. The result does not
assert that $D_r$ is attained or that the number exceeds $K$.
It is neither an OOD prediction lower bound nor a simultaneous
dimension-width theorem.

\section{Critical law and its mean}
Write $U=(b,w)$ and $G(U)=\E_x\chi(x)\relu(U\cdot(1,x))$.
The repository's sharp correlation inequality is
\[
 |G(U)|\le\rho\sqrt r\min_i|w_i|\le\rho\|U\|.
\]
Equality in the last inequality occurs exactly at
$v_z=(0,z/\sqrt r)$, $z\in X$, up to positive scaling.
If $\kappa=G(v_{\one})$, then $|\kappa|=\rho$ and the critical
parameter rays are
\[
 \mathcal R=\{R(\sigma_z,v_z):R\ge0,\ z\in X\},\qquad
 \sigma_z=\operatorname{sgn}(\kappa)\chi(z).
\]
The continuous homogeneous quadratic function
\[
 D_0(\theta)=\tfrac12\|\theta\|^2-Z(\theta)/\rho
\]
is nonnegative and vanishes precisely on $\mathcal R$.
The density of $Q_*$ is a fixed positive constant times $e^{-D_0}$.

For completeness, endpoint fourth-moment finiteness follows from the
sharp corner calculation in the companion note. In a cone about a
maximizing incoming ray, write $U=Rv_z+u$, $u\perp v_z$. Then
$|G(U)|=\rho R-h_z(u)$, where $h_z$ is a norm on the $r$-dimensional
transverse space. After integrating out the readout, the exponent is
comparable to $-R\|u\|$. The substitution $u=v/R$ shows that a
fourth parameter moment is bounded by a multiple of
$\int_1^\infty R^{4-r}dR$, finite for $r>5$.
Outside the cones there is a strict Gaussian quadratic gap.

The analytic bound in the softplus companion applies without change
to ReLU. With $T=\min_i|w_i|$ and $c=\sqrt{2/\pi}$,
\[
 m_c\le\rho r^2c\int_0^\infty t^2
       [e^{t^2/2}2\overline\Phi(t)]^{r-1}dt
 \le\frac{\pi\rho r^2}{(r-2)(r-3)(r-4)}<\frac38.           \tag{4}
\]
The first inequality integrates the Gaussian readout and uses
$M_*\ge1$. The second uses
$e^{t^2/2}2\overline\Phi(t)\le(1+ct)^{-1}$, proved by the elementary
one-sign-change differential comparison in that note. The final bound
holds at $r=6$, where it is $9\pi/(32\sqrt6)$, and decreases over
subsequent even dimensions. Cube sign symmetries force
$\E_{Q_*}\Phi=m_c\chi$.

\section{A polynomial lower bound for a centered small box}
We give the local probability argument explicitly, so no unverified
local central-limit assumption enters the conditioning step.

\begin{lemma}[A Gaussian component gives a local lower bound]\label{lem:local}
Let $V\in\R^m$ have finite second moment and mean $\mu$.
Suppose some convolution power of its law has a centered nondegenerate
Gaussian component: for some integer $b\ge1$, $p>0$, and $\sigma>0$,
\[
 \mathcal L(V_1+\cdots+V_b)
       =pN(0,\sigma^2I_m)+(1-p)R.
\]
Then, for every fixed $a>0$ and all sufficiently large $n$,
\[
 \Pp\{\|V_1+\cdots+V_n-n\mu\|_\infty<a\}
              \ge c_a n^{-m/2},
\]
where $c_a>0$ may depend on the law and the dimension.
\end{lemma}
\begin{proof}
First take $n=bq$. Represent each block by an independent Bernoulli
choice of the Gaussian component or $R$. Let $J$ count the Gaussian
blocks and let $T$ sum the non-Gaussian blocks, treating Gaussian
choices as zero in $T$. Then $\E T=n\mu$ and
$\E\|T-n\mu\|^2\le Cq$. The residual law has finite second moment
because the full block law and the Gaussian component do.
Chebyshev and the binomial lower-tail bound give a fixed positive
probability to
\[
 \|T-n\mu\|\le C'\sqrt q,\qquad pq/2\le J\le q.
\]
Conditional on these choices and residual values, the remaining sum
is $N(0,J\sigma^2I_m)$. Throughout a fixed box about $n\mu-T$ its
density is bounded below by $cq^{-m/2}$, uniformly on this event.
Integrating the box proves the claim. If $n$ is not a multiple of
$b$, condition the at most $b-1$ remaining variables to a fixed bounded
set of positive probability; their bounded offset changes no estimate.
\end{proof}

The lemma applies to $V=\Phi(\theta)$ under $Q_*$. ReLU features on
the finite cube span $\R^M$: use a linear score taking distinct values,
a constant positive feature, and hinges between consecutive score
values, obtaining an invertible triangular evaluation matrix.
Choose $M$ incoming vectors in compact neighborhoods of such a basis.
Under $Q_*$, each incoming law has a strictly positive density.
Conditional on the incoming vectors, the readouts are independent
Gaussians with variance one and means $G(U)/\rho$.
Their output sum is therefore a full-dimensional Gaussian whose mean
is uniformly bounded and whose covariance lies between two fixed
positive multiples of the identity on those neighborhoods.
Every Gaussian in this compact family has density at least a fixed
positive multiple of a sufficiently narrow centered Gaussian density,
uniformly on all of $\R^M$. Multiplying by the positive probability
of the incoming neighborhoods gives the required Gaussian component
of the $M$-fold output convolution. Fourth parameter moments imply
second output moments. Thus Lemma~\ref{lem:local} applies with $m=M$.

\section{Polynomial evidence and posterior domination}
Put $\delta=s-m_c>0$ and use the canonical hemisphere
$\mathcal S=\{z:z_1=1\}$, of size $K$. The repository's Fourier
identity gives
\[
 \sum_{z\in\mathcal S}\sigma_z\relu(z\cdot x/\sqrt r)
                         =K\rho\chi(x),\qquad r\ge4.       \tag{5}
\]
Reserve $K$ fixed neuron indices at centers
$R_N(\sigma_z,v_z)$, where
\[
 R_N^2=\frac{N\delta+Km_c+2}{K\rho}.
\]
Their total output is $(N\delta+Km_c+2)\chi$.
Constrain the other $N-K$ neurons' total output to a fixed small box
about $(N-K)m_c\chi$. By the preceding lemma this has probability
at least $cN^{-M/2}=cN^{-K}$ under $Q_*^{N-K}$.

Constrain each reserved parameter to a ball of radius $c'/\sqrt N$
about its center. Choose $c'>0$ small enough that the combined output
perturbation is at most $1/4$ at each input. This is possible because
ReLU output has local Lipschitz constant $O_r(\sqrt N)$ there.
At the centers $D_0=0$ and throughout these balls $D_0=O_r(1)$.
Their total $Q_*$ probability is at least
$cN^{-K(r+2)/2}$. On the combined event all signed total outputs
are between $Ns+1$ and $Ns+3$.

The exact change of measure is
\[
 P_0^N(H_N)=e^{-NI(s)}
 \E_{Q_*^N}\left[
 e^{-t_*(\sum_jZ_j-Ns)}\one_{H_N}\right].                  \tag{6}
\]
The exponential factor is at most one on $H_N$, and is bounded below
by a positive constant on the constructed event. This proves (1).
It also proves the useful domination, for every event $A$,
\[
 P_0^N(A\mid H_N)\le C_r N^{D_r}Q_*^N(A).                 \tag{7}
\]
Bounded-test concentration under the product law, followed by (7),
proves weak empirical convergence to $Q_*$. A countable convergence-
determining family plus tightness gives the probability-measure statement.
For tightness, choose a compact set with arbitrarily large $Q_*$ mass
and apply the same Bernoulli concentration to its empirical mass.

On $H_N$, if one signed average output exceeds $s+u$, then
$N^{-1}\sum_jZ_j\ge s+u/M$. Equation (6) and its polynomial lower
bound make its posterior probability at most
$C_rN^{D_r}e^{-Nt_*u/M}$. Hence $F_N\to s\chi$.

\section{Only finitely many large neurons can carry the missing mean}
First, total energy is bounded. Under the original Gaussian prior,
\[
 P_0^N\left\{N^{-1}\sum_j\|\theta_j\|^2>B\right\}
 \le\exp\{-N[B/4-(r+2)\log2/2]\}.
\]
Choose $B>4[I(s)+(r+2)\log2/2+1]$ and divide by the evidence lower
bound. The posterior probability tends to zero exponentially.

For any fixed $\varepsilon>0$, fourth-moment Markov gives
\[
 Q_*\{\|\theta\|>\varepsilon\sqrt N\}\le C_\varepsilon N^{-2}.
\]
Independence and a union bound over index sets show that the chance of
$D_r+1$ or more such parameters is $O_\varepsilon(N^{-D_r-1})$.
By (7), the posterior count is at most $D_r$ with probability tending
to one.

For any fixed angular separation from $\mathcal R$, the function
$D_0$ has a strictly positive minimum on the corresponding part of
the unit sphere. On the annulus
$\varepsilon\sqrt N\le\|\theta\|\le\sqrt{BN}$, its $Q_*$ density
is therefore at most $Ce^{-cN}$. The annulus volume is polynomial
in $N$. Union over the neurons, use (7), and use the energy bound
for the outside region. Thus all these large parameters approach
$\mathcal R$ in angle in posterior probability.

We also need to exclude a collective signal from arbitrarily many
smaller neurons. There is a constant $C_r$ such that
$\|\Phi(\theta)\|_\infty\le C_r\|\theta\|^2$.
Truncate at $\|\theta\|\le\varepsilon\sqrt N$ and denote the
truncated output vector by $V^{(N)}$. Its coordinate second moments
are uniformly bounded, and its mean tends to $m_c\chi$.
For $\lambda=A\log N/N$, the elementary exponential Taylor bound gives
\[
 \log\E_{Q_*}e^{\lambda(V_x^{(N)}-\E V_x^{(N)})}
 \le C\lambda^2\exp(2\lambda C_r\varepsilon^2N)
\]
for large $N$; the harmless factor two also covers centering.
Choose $\varepsilon$ so small that $2AC_r\varepsilon^2<1/2$.
Multiplying the log moment by $N$ makes its upper bound $o(1)$.
The same estimate for $-V_x^{(N)}$ and Markov's inequality imply,
for any fixed $u>0$,
\[
 Q_*^N\left\{
 \left\|N^{-1}\sum_jV_j^{(N)}-m_c\chi\right\|_\infty>u
 \right\}\le 2M N^{-Au/2+o(1)}.                           \tag{8}
\]
For any desired polynomial power first choose $A$ large and then
$\varepsilon$ small. In particular this probability is $o(N^{-D_r})$,
so the same truncated output is within $u$ of $m_c\chi$ in posterior
probability. No exponential output moment under $Q_*$ was assumed.

\section{The parity geometry forces exponentially many large neurons}
A finite collection on the critical rays contributes
\[
 \sum_{z\in X} m_z\sigma_z h_z(x),\qquad
 h_z(x)=\relu(z\cdot x/\sqrt r),\quad m_z\ge0.
\]
The repository's nonvanishing even-Fourier-coefficient calculation
implies that equality to $\delta\chi$ forces
\[
 m_z+m_{-z}=\frac{\delta}{K\rho}>0
                          \quad\text{for every antipodal pair}. \tag{9}
\]
Every pair must therefore contain a nonzero neuron. In particular
fewer than $K$ critical-ray neurons cannot supply $\delta\chi$.

For a quantitative separation sufficient here, restrict the normalized
parameter energy to the fixed bound $B$. The output images of all
collections of at most $K-1$ critical-ray parameters in this bounded
set form a compact set: there are finitely many ray choices, and zero
parameters may pad the collection. By (9) this compact set excludes
$\delta\chi$, so their sup-norm distance from it is some $g_r>0$.
Choose $u<g_r/4$, and then choose $\varepsilon_r$ small enough for
(8) to beat the domination power $D_r$.

With posterior probability tending to one, total output is within $u$
of $s\chi$, the truncated small-neuron output is within $u$ of
$m_c\chi$, total energy is at most $B$, and all parameters above
the cutoff are arbitrarily close to critical rays. If there were
fewer than $K$ such parameters, their summed output would be within
$2u+o(1)<g_r$ of $\delta\chi$, contradicting the compact separation.
This proves the lower count in (2); the earlier fourth-moment argument
proves its upper count.

\section{Energy accounting and the learned feature kernel}
The same truncation argument applies to any finite collection of
functions bounded by a constant times $\|\theta\|^2$, since fourth
parameter moments are finite. In particular it applies to $a^2$, to
$\|U\|^2$, and to each product
$\relu(U\cdot\bar x)\relu(U\cdot\bar y)$.
For parameters on the critical rays,
\[
 a^2=\|U\|^2=Z/\rho.
\]
The angular convergence and the total-energy bound make the combined
error in these identities tend to zero for all parameters above any
fixed cutoff. The small parameters contribute arbitrarily close to
their $Q_*$ expectations by (8). Since the total parity coefficient
tends to $s$, it follows that the readout and incoming energy limits
are their respective critical expectations plus $(s-m_c)/\rho$.

Conditional Gaussian readouts give
$\E_{Q_*}a^2=1+m_c/\rho$. The incoming tilted density is proportional
to $\exp[-(\|U\|^2-G(U)^2/\rho^2)/2]$.
Its exponent is homogeneous of degree two; scaling its normalizing
integral, or polar integration, gives
$\E_{Q_*}[\|U\|^2-G(U)^2/\rho^2]=r+1$.
Since $\E_{Q_*}G(U)^2/\rho^2=m_c/\rho$, this proves (3).

There is also a deterministic statement about the even Fourier block
of the learned feature kernel. Define
\[
 \mathcal K_N(x,y)=N^{-1}\sum_j
       \relu(U_j\cdot\bar x)\relu(U_j\cdot\bar y),\qquad
 \mathcal K_*(x,y)=\E_{Q_*}
       [\relu(U\cdot\bar x)\relu(U\cdot\bar y)].
\]
Let $h(x)=\relu(\sum_i x_i/\sqrt r)$ and let $T,V\subseteq[r]$
have even cardinality, including zero. Then in posterior probability,
\[
 \E_{x,y}\chi_T(x)[\mathcal K_N-\mathcal K_*](x,y)\chi_V(y)
 \longrightarrow
 \frac{s-m_c}{\rho}\widehat h(T)^2\one_{T=V}.              \tag{10}
\]
Indeed the small neurons give $\mathcal K_*$ by the same truncated
moment argument. For the large critical-ray parameters, the output
increment determines their antipodal-pair masses by (9), up to errors
that vanish as the truncation tolerance vanishes. On even characters,
opposite orientations have identical feature coefficients. Summing
the fixed pair masses and using character orthogonality on the
antipodal quotient gives (10). Every even coefficient of $h$ is
nonzero by the existing repository lemma, so the extra even kernel
block has rank $K=2^{r-1}$.

Equation (10) is a rank statement, not an effective-rank or numerical-
conditioning claim. Some eigenvalues can be very small. It shows that
weak convergence of ordinary-neuron laws alone misses part of both
the fitted function and its feature kernel.

\section{An exponential increase in the target kernel eigenvalue}
The kernel calculation has a stronger interpretation than its rank alone.
For any kernel $\mathcal K$ define its parity eigenvalue by
\[
 \lambda_\chi(\mathcal K)=\E_{x,y}\chi(x)\mathcal K(x,y)\chi(y).
\]
For the invariant kernels $\mathcal K_0$ and $\mathcal K_*$ this is
the eigenvalue of the parity character, with uniform measure in the
kernel operator. For a finite empirical kernel it denotes its quadratic
form in that direction; the limiting even Fourier block is diagonal.

\begin{theorem}[Target feature amplification]
For the ordinary unit-Gaussian ReLU prior kernel $\mathcal K_0$,
\[
 0<\lambda_\chi(\mathcal K_0)\le3\,2^{-r},\qquad r\ge6\text{ even}.
\]
In the actual posterior of the first theorem,
\[
 \mathcal K_N\chi\longrightarrow s\rho\,\chi,\qquad
 \lambda_\chi(\mathcal K_N)\longrightarrow s\rho
 \ge \frac{s}{2\sqrt2\,r}.
\]
In particular the target feature eigenvalue grows by a factor at least
$s2^r/(6\sqrt2\,r)$. The exceptional neurons alone supply
$(s-m_c)\rho>\rho/8$ at $s=1/2$.
\end{theorem}
\begin{proof}
The critical ordinary law has
\[
 \lambda_\chi(\mathcal K_*)=\E_{Q_*}G(U)^2=\rho m_c,
\]
because the conditional readout mean is $G(U)/\rho$. Taking $T=V=[r]$
in (10) adds $(s-m_c)\rho$, proving the quadratic-form formula.
For the full eigenvector assertion, on a critical ray
$G(U)\relu(U\cdot\bar x)=\rho\Phi_x(\theta)$.
The truncated quadratic-moment argument therefore makes the large-neuron
part of $\mathcal K_N\chi$ tend to $\rho(s-m_c)\chi$; the small
part tends to $\mathcal K_*\chi=\rho m_c\chi$.
The elementary
central-binomial estimate in the existing parity calculation gives
$\rho\ge1/(2\sqrt2\,r)$.

For the prior bound use the classical Gaussian ReLU kernel formula
(\href{https://proceedings.neurips.cc/paper/2009/file/5751ec3e9a4feab575962e78e006250d-Paper.pdf}{Cho and Saul, 2009}). Its two preactivations
have variance $r+1$ and correlation $c=(1+x\cdot y)/(r+1)$, and
\[
 \mathcal K_0(x,y)=(r+1)\left\{\frac c4+
 \frac1{2\pi}\left[1+\sum_{l\ge1}a_lc^{2l}\right]\right\},\qquad
 a_l=\frac{\binom{2l}{l}}{4^l(2l-1)^2}\le l^{-5/2}.
\]
The series is absolutely convergent even at the endpoints. The last
inequality follows from
$\binom{2l}{l}/4^l\le(l+1)^{-1/2}$, proved by induction, and
$2l-1\ge l$.
For uniform independent $\varepsilon_i\in\{-1,1\}$ put
\[
 A_n=\E\left[\prod_i\varepsilon_i
      \left(\frac{1+\sum_i\varepsilon_i}{r+1}\right)^n\right].
\]
Expanding the power identifies $A_n$ with the probability that, in $n$
independent uniform draws from $\{0,1,\ldots,r\}$, every nonzero symbol
has odd multiplicity. Thus $A_n=0$ for $n<r$ and $A_n\ge0$.
Poissonize the counts: independent $P_i\sim\mathrm{Poisson}(n/(r+1))$
conditioned on $\sum_{i=0}^rP_i=n$ have those multinomial counts. Hence
\[
 A_n\le\frac{\prod_{i=1}^r\Pp(P_i\text{ odd})}
                    {\Pp(\mathrm{Poisson}(n)=n)}
     \le e\sqrt n\,2^{-r}.
\]
For the denominator, the elementary factorial bound
$n!\le e\sqrt n(n/e)^n$ follows by integrating the concave function
$\log x$ over $[1,n]$ and comparing with its trapezoids.
The constant and linear kernel terms have zero parity coefficient.
Consequently
\begin{align*}
 \lambda_\chi(\mathcal K_0)
 &\le\frac{e\sqrt2(r+1)}{2\pi}\,2^{-r}
                \sum_{l\ge r/2}l^{-2}\\
 &\le\frac{e\sqrt2(r+1)}{2\pi(r/2-1)}\,2^{-r}
 <3\,2^{-r}.
\end{align*}
Positivity follows already from the $n=r$ term, where every nonzero
symbol can occur exactly once. This proves the claim.
\end{proof}

This is feature adaptation in a directly measurable kernel direction:
the initial parity eigenvalue is exponentially small, while its learned
value is of order $1/r$. It does not supply a sample-complexity or OOD
claim, since this experiment conditions on every input. Nor does the
exponential exceptional-neuron count give a width lower bound at fixed
relative energy accuracy: the separate Monte Carlo construction achieves
$(1+\epsilon)$ of optimal margin energy using polynomial width.

\section{Almost all ordinary neurons remain nearly Gaussian at large parity degree}
The preceding prior-kernel bound also makes the critical mean exponentially
small in the parity dimension, improving the earlier uniform bound $m_c<3/8$.
Put $c=\sqrt{2/\pi}$ and $u_0=1+2c/3$. Then
\[
 m_c\le \frac3\rho\exp[-r(\log2-2/9)]
       +\frac{\rho r^2}{c^2(r-4)}u_0^{-(r-4)}.            \tag{13}
\]
In particular $m_c=O(r e^{-c_0r})$, where $c_0=\log u_0>0.42$.
Consequently
\[
 \operatorname{KL}(Q_*\|P_0)\le m_c/\rho
       =O(r^2e^{-c_0r}),\qquad
 \operatorname{TV}(Q_*,P_0)=O(r e^{-c_0r/2}).              \tag{14}
\]
These statements compare one-neuron probability laws. They do not imply
small distance between the full conditioned network and its prior.

To prove (13), write $T=\min_i|w_i|$, and integrate out the readout:
\[
 m_c=\frac1{\rho M_*}\E_0[G(U)^2e^{G(U)^2/(2\rho^2)}].
\]
On $T\le a$, use $|G(U)|\le\rho\sqrt r\,T$ and
$\E_0G(U)^2\le3\,2^{-r}$ to bound this contribution by
$(3/\rho)e^{-r(\log2-a^2/2)}$.
On $T>a$, the same sharp inequality and the density of the minimum of
$r$ independent half normals give the upper bound
\[
 \rho r^2c\int_a^\infty t^2
       [e^{t^2/2}2\overline\Phi(t)]^{r-1}\,dt
 \le \frac{\rho r^2}{c^2(r-4)}(1+ca)^{-(r-4)}.
\]
The last step uses the Mills-ratio bound already proved above and
$t^2\le(1+ct)^2/c^2$. Taking $a=2/3$ proves (13).
The central-binomial estimates give $\rho=\Theta(1/r)$, and
$\log2-2/9>\log u_0$, proving the stated exponential rate.
Finally, $\operatorname{KL}(Q_*\|P_0)=m_c/\rho-\log M_*$,
$M_*\ge1$, and Pinsker's inequality give (14).

In the iterated limit (first width, then even parity dimension), the
ordinary-neuron law therefore approaches the original Gaussian law.
Nonetheless the exceptional neurons carry asymptotically all of the
parity output and parity-kernel eigenvalue. This is a precise sense in
which weak or bounded-observable descriptions of the typical neuron
can miss essentially the entire learned feature.

\section{The same transition under square or logistic likelihood}
The effect does not require a hard margin indicator. Keep the same
prior and put a width-scaled likelihood on all parity examples:
\[
 \frac{d\Pi_{N,\lambda}}{dP_0^N}
 \ \propto\ \exp\left[-\frac{N\lambda}{M}
                \sum_{x\in X}\ell(\chi(x)F_N(x))\right],
 \qquad \lambda>0.
\]
We consider either square loss $\ell(s)=(1-s)^2/2$ or logistic loss
$\ell(s)=\log(1+e^{-s})$. The factor $N$ is part of the model: this
is the mean-field, width-scaled likelihood, not a fixed-noise likelihood
on a fixed dataset as width diverges.
Let $\Lambda(t)=\log\E_0e^{tZ}$, $Q_t=e^{tZ-\Lambda(t)}P_0$,
and $m(t)=\Lambda'(t)$ for $0\le t\le t_*$.

\begin{theorem}[A finite likelihood-strength transition]
Define
\[
 \lambda_c=
 \begin{cases}
 t_*/(1-m_c),&\text{square loss},\\
 t_*(1+e^{m_c}),&\text{logistic loss}.
 \end{cases}
\]
For $0<\lambda<\lambda_c$, there is a unique $t\in(0,t_*)$
solving $t=-\lambda\ell'(m(t))$. The posterior empirical law tends
to $Q_t$, the network output tends to $m(t)\chi$, and the largest
parameter norm divided by $\sqrt N$ tends to zero.

At $\lambda=\lambda_c$, the same statements hold with $t=t_*$
and $m(t)=m_c$, still with no macroscopic parameters.

For $\lambda>\lambda_c$, the ordinary-neuron law remains $Q_*$,
but the output tends to $s_\lambda\chi$, where
\[
 s_\lambda=
 \begin{cases}
 1-t_*/\lambda,&\text{square loss},\\
 \log(\lambda/t_*-1),&\text{logistic loss}.
 \end{cases}
\]
There is an $\varepsilon_{r,\lambda}>0$ for which the exceptional-
neuron count is between $K$ and $D_r$, with probability tending to
one. The energy limits (3) and the even-kernel formula (10) hold with
$s=s_\lambda$. If $\mathcal Z_{N,\lambda}$ is the likelihood evidence,
\[
 -\log\mathcal Z_{N,\lambda}
   =N[I(s_\lambda)+\lambda\ell(s_\lambda)]+O_{r,\lambda}(\log N).
\]
In particular the entirely explicit choice $\lambda=2/\rho$ under
square loss gives $s_\lambda=1/2$ and the exponential-neuron conclusion.
\end{theorem}
\begin{proof}
The function $m(t)$ is continuous and strictly increasing, starts at
zero, and tends to $m_c<1$ at $t_*$. For either loss,
$t+\lambda\ell'(m(t))$ is strictly increasing in $t$, is negative
at zero, and is positive at $t_*$ precisely when
$\lambda<\lambda_c$. This proves the unique subcritical solution.

Write $y_x=\chi(x)F_N(x)$. Changing measure to $Q_t^N$ leaves the
potential
\[
 V_t(y)=\frac1M\sum_x[t y_x+\lambda\ell(y_x)].
\]
For the selected $t$ this strictly convex, coercive function has its
unique minimum at $m(t)\one$. Under $Q_t^N$ the empirical outputs
are within $C/\sqrt N$ of that point with probability bounded below,
by their finite second moments and Chebyshev's inequality. A local
quadratic upper bound on $V_t-\min V_t$ therefore bounds the
normalizing expectation of $e^{-N(V_t-\min V_t)}$ below by a fixed
positive constant. The posterior is dominated by a constant times
$Q_t^N$. This proves output and empirical-law convergence and absence
of macroscopic parameters, using Gaussian tails for $t<t_*$ and the
finite fourth parameter moment at $t=t_*$. It also proves that the
critical case has no macroscopic neurons.

Above the transition, change measure to $Q_*^N$. Now $V_{t_*}$ has
its unique minimum at $s_\lambda\one$, as the displayed formulas
show, and $s_\lambda>m_c$. The polynomial local construction used
for (1), with centers supplying $(s_\lambda-m_c)\chi$, places the
output within $O(1/N)$ of this minimum at $Q_*^N$ probability
at least $cN^{-D_r}$. Its likelihood factor is bounded below there.
The factor is at most one everywhere after subtracting the minimum,
so the posterior is dominated by $CN^{D_r}Q_*^N$ and the evidence
has the stated asymptotics. Coercivity gives output concentration.
The original-prior Gaussian energy bound still applies because both
losses are nonnegative. All truncation, ray, count, energy, and kernel
arguments above consequently apply without change.
\end{proof}

\section{Two exact finite-width identities}
The asymptotic energy and kernel formulas have useful finite-width
checks. Let $L(F)$ be any bounded nonnegative likelihood with positive
evidence, including the hard margin indicator, under the unit-Gaussian
prior. Positive homogeneity makes each neuron's output invariant under
$a_j\mapsto e^ta_j$, $U_j\mapsto e^{-t}U_j$. Changing variables in the
Gaussian integral and differentiating at zero gives, exactly,
\[
 \E_{\Pi}[\|U_j\|^2-a_j^2]=r,\qquad
 \E_{\Pi}\left[N^{-1}\sum_j(\|U_j\|^2-a_j^2)\right]=r.       \tag{11}
\]
Indeed the forward change-of-variable Jacobian is $e^{-rt}$,
whereas the derivative of the Gaussian exponent after the substitution
is $\|U_j\|^2-a_j^2$. Their sum has expectation zero.
Gaussian domination justifies differentiation,
without differentiating the possibly discontinuous likelihood.
For unequal layer variances replace the squares by their standardized
squares. This identity does not require large width or parity labels.

For the square likelihood $L(F)=\exp[-N\lambda\|F-f\|_2^2/2]$,
with the uniform input norm, Gaussian integration by parts in each
readout gives another exact identity:
\[
 \E_\Pi F=\lambda\,\E_\Pi[\mathcal K_N(f-F)].               \tag{12}
\]
Here $\mathcal K_N$ acts by uniform averaging over its second argument.
The features do not depend on the readouts, so no derivative of an
activation is needed. In a concentrating parity phase $F\to s\chi$,
this is consistent with the target kernel eigenvalue $s/t$, where
$t=\lambda(1-s)$. The critical value $t=t_*$ gives $s\rho$ as above.
These are checks on the posterior, not convergence guarantees for a
particular sampler or optimizer.

\section{An independent sampler for the critical ordinary law}
The critical-law predictions can be checked without a neural posterior
Markov chain. Here is an exact rejection envelope, in real arithmetic,
for the incoming marginal of $Q_*$ at even $r>2$.
Put $c=\sqrt{2/\pi}$ and
\[
 A(t)=\int_0^\infty e^{-tv-v^2/2}dv
     =\sqrt{2\pi}e^{t^2/2}\overline\Phi(t).
\]
The earlier Mills-ratio inequality gives
$A(t)\le c^{-1}(1+ct)^{-1}$. Generate:
\begin{enumerate}
\item $t\ge0$ with density $(r-2)c(1+ct)^{-(r-1)}$;
      accept this first stage with probability
      $[cA(t)(1+ct)]^{r-1}$.
\item Pick a minimum coordinate uniformly. Give it magnitude $t$.
      Independently sample the other excess magnitudes $v_i\ge0$ with
      density $e^{-tv_i-v_i^2/2}/A(t)$, and set their magnitudes to
      $t+v_i$. Give the weights independent uniform signs and take
      $b\sim N(0,1)$.
\item Accept with probability
      $\exp[-(rt^2-G(b,w)^2/\rho^2)/2]\le1$.
      The inequality follows from $|G|\le\rho\sqrt r\,t$.
\end{enumerate}
After acceptance take $a\sim N(G(b,w)/\rho,1)$ independently
conditional on the incoming parameters. The resulting whole neuron has
law $Q_*$. In permutation- and sign-invariant observables such as $G^2$
and incoming energy, one may fix the minimum coordinate and all weight
signs to save computation.

To verify the algorithm, fix the positive orthant and its unique minimum
coordinate, and write $w_i=t+v_i$ with one $v_i=0$. Multiplying the
proposal density by both acceptance probabilities cancels every $A(t)$
and $1+ct$ factor, and gives
\[
 (r-2)c^r\phi(b)
 \exp\left[-\frac12\sum_{i\ne\min}v_i^2-t\sum_{i\ne\min}v_i
             -\frac{rt^2}{2}+\frac{G(b,w)^2}{2\rho^2}\right].
\]
The unnormalized tilted Gaussian incoming density in these coordinates,
after summing its $r$ possible minimum positions, is the same expression
with leading coefficient $rc^r$. Thus the total acceptance probability is
\[
 p_{\rm acc}=\frac{r-2}{r}M_*,\qquad
 1\le M_*\le\frac r{r-2}.
\]
The acceptance probability is at least $(r-2)/r$, hence at least $2/3$
for the posterior theorem's dimensions. This is a sampler for the
\emph{limiting ordinary neuron}, not for the full finite-width conditioned
network or its exceptional-neuron orientations. The coordinate signs
and readout draw must be restored when sampling noninvariant observables.
The implementation uses floating-point arithmetic and is not an interval
certificate. For $r\ge6$ the finite fourth parameter moment gives finite
variance to the Rao--Blackwell estimator $G(U)^2/\rho$ of $m_c$.

\section{Small teacher, scaling variants, and scope}
On the grid $S(x)=\sum_i x_i\in\{-r,-r+2,\ldots,r\}$,
\[
 \chi(x)=1-\relu(S+r)
       +2\sum_{j=1}^{r-1}(-1)^{j+1}\relu(S+r-2j).
\]
Implement the constant one by $\relu(1)$. This is a width-$r+1$
teacher with exact signed output one, rational weights, and magnitudes
at most $r$ in the sum convention. Multiplying all readouts by $r+1$
implements the averaged convention with weight bound $2r+2$.

The sharper companion proves exactly $K$ ReLU exceptional neurons,
one per antipodal pair with equal limiting masses, using local polynomial
volume estimates; the orientation-law companion identifies their limiting
pattern probabilities. The softplus companion obtains the same count using a
second exponential order of speed $\sqrt N$; ReLU instead has
polynomial critical tails. No positive or negative
prediction claim follows from either representation theorem by itself.
The present argument is conventional, not Lean verified or independently
reviewed.
\section{Gaussian normalization does not remove the mechanism}
For independent Gaussian layer scales $a=\tau_a\widetilde a$ and
$U=\tau_u\widetilde U$, positive homogeneity identifies the
physical-margin-$s$ posterior exactly with the unit-scale posterior at
$\widetilde s=s/(\tau_a\tau_u)$. The critical margin is therefore
$\tau_a\tau_u m_c$, and the exceptional count is still $K$ whenever
$s$ exceeds this value. The learned incoming-kernel parity eigenvalue
becomes $s\rho\tau_u/\tau_a$, because that kernel scales by
$\tau_u^2$. The total layer energies become
\[
 \frac1N\sum_j a_j^2\longrightarrow
       \tau_a^2+\frac{s\tau_a}{\rho\tau_u},\qquad
 \frac1N\sum_j\|U_j\|^2\longrightarrow
       (r+1)\tau_u^2+\frac{s\tau_u}{\rho\tau_a}.
\]
For fan-in scaling $\tau_u=r^{-1/2}$, $\tau_a=1$, the learned
eigenvalue is of order $r^{-3/2}$ rather than $r^{-1}$, still
exponentially larger than its initial parity eigenvalue. These exact
rescaling statements fix the dimension and positive layer scales before
taking width to infinity; they add no uniform bound on the required
width or a cold-prior interchange of limits.
\end{document}
